同一个变换，可以写成两个长得毫无关系的矩阵。

取平面上``沿某条斜线方向拉伸两倍''这个操作。用标准坐标记录，它是四个格子都非零的矩阵；把坐标轴转到与那条斜线对齐，它就变成对角阵 \(\diag(2,1)\)。两个矩阵没有一个数字相同，做的却是同一件事。数字属于记账方式，操作本身与坐标无关。

这个观察不是修辞，它是整个线性代数的组织原则。矩阵里那些纠缠在一起的数，很多时候只是坐标没选好留下的痕迹。于是这门学科反复在做同一个动作：\textbf{给手上的问题挑一组基，让原本互相牵扯的操作变成互不干扰的一维伸缩。}分解定理之所以一个接一个，是因为``什么样的问题配什么样的基''有不同的答案。

\begin{center}
\begin{tikzpicture}[x=1cm,y=1cm,font=\small,>=stealth]
  \draw[black!45,fill=blue!9] (0,-0.75) rectangle (3.0,0.75);
  \node[align=center] at (1.5,0) {挑一组\\合适的基};
  \draw[black!45,fill=green!8] (4.6,2.0)  rectangle (13.0,3.1);
  \node at (8.8,2.55) {正交基：投影退化成逐项内积};
  \draw[black!45,fill=green!8] (4.6,0.3)  rectangle (13.0,1.4);
  \node at (8.8,0.85) {特征基：反复作用退化成标量各自取幂};
  \draw[black!45,fill=green!8] (4.6,-1.4) rectangle (13.0,-0.3);
  \node at (8.8,-0.85) {奇异向量基：任意矩阵退化成沿轴伸缩};
  \draw[black!45,fill=green!8] (4.6,-3.1) rectangle (13.0,-2.0);
  \node at (8.8,-2.55) {Fourier 基：卷积退化成逐点相乘};
  \draw[->,black!55] (3.05,0.25) -- (4.55,2.55);
  \draw[->,black!55] (3.05,0.12) -- (4.55,0.85);
  \draw[->,black!55] (3.05,-0.12) -- (4.55,-0.85);
  \draw[->,black!55] (3.05,-0.25) -- (4.55,-2.55);
\end{tikzpicture}

\emph{九章看上去在讲九种分解，动作只有一个：换坐标，把纠缠的操作拆成互不干扰的一维动作。区别只在于``什么算合适''。}
\end{center}

\subsection{五条贯穿始终的判断}

\textbf{矩阵是线性变换的坐标记录，不是数字表格。}把 \(Ax=b\) 读成``\(b\) 能不能由 \(A\) 的列向量拼出来''，比读成``解一组方程''有用得多——前者立刻回答有没有解、解唯一不唯一，也立刻解释了为什么回归里的 \(Ax=b\) 几乎总是无解，只能退而求投影。同一个变换换基后数字全变，维数、秩、特征值这些量却不动。能不能分清哪些属于操作本身、哪些只是记账方式，是读懂这一部分的门槛。

\textbf{正交把``最近''变成可计算的。}残差必须垂直于列空间，这条几何事实直接生出正规方程，最小二乘因此不是拟合技巧而是投影。同一句话还解释了拟合之后残差为什么与每个特征都不相关——若还相关，说明你没投到底。

\textbf{对称性换来谱定理，正定性接上优化。}对称矩阵一定有一组正交特征基，于是二次型的几何形状完全由特征值符号决定：全正是碗，有正有负是鞍。协方差矩阵天然对称半正定，Hessian 在极小点附近正定，条件数 \(\lambda_{\max}/\lambda_{\min}\) 直接决定梯度下降是顺畅下滑还是在峡谷里来回撞墙。这条线把线性代数与后面的优化焊在一起。

\textbf{SVD 没有任何前提，这是它成为数据工具的原因。}特征分解要求方阵且可对角化，真实的数据矩阵两个条件一个都不满足。SVD 对任意矩阵都成立，奇异值有序排列，于是``保留前 \(k\) 个''就是最优低秩近似。降维、压缩、去噪、矩阵补全、低秩适配共用的是同一条定理；而队尾那些接近零的奇异值正是病态的来源——正则化做的事，说白了就是不去除以接近零的数。

\textbf{好性质到高阶会大面积失效。}三阶以上的张量没有 SVD 的直接推广，最优低秩近似甚至可能根本不存在。这不是技术遗憾，它解释了为什么张量分解只能给局部解、结果依赖初始化，与矩阵 SVD 的确定性恰成对照。把``矩阵的经验能不能往上推''当作一个需要验证的问题，而不是默认的延伸。

\subsection{九章怎么安排}

\hyperref[ch:017]{``线性方程组与矩阵''}和\hyperref[ch:018]{``向量空间、基与线性变换''}是地基，必须扎实。前者建立行视角与列视角的对照，后者把四个基本子空间摊开——方程何时有解、解是否唯一、冗余藏在哪里，三个问题在这里一次回答完。

\hyperref[ch:019]{``正交、投影与最小二乘''}是第一个回报点，也是统计与机器学习用得最多的一章。\hyperref[ch:021]{``特征值、对角化与动力系统''}和\hyperref[ch:022]{``对称矩阵、二次型与正定性''}应当连读：前者管一般方阵，后者说明``对称''这个假设能额外买到什么。

\hyperref[ch:023]{``SVD：从任意矩阵到数据中的主要方向''}是全部分的枢纽。时间紧张的话，读完地基两章后可以直接跳到这里，再回头补特征值——顺序反过来更顺，但先看到 SVD 能让人明白前面在铺什么。

\hyperref[ch:020]{``行列式：一个数怎样概括线性变换''}篇幅短，按体积伸缩因子理解即可，不必在余子式展开上停留；它真正的用处出现在概率密度的变量替换里。\hyperref[ch:024]{``复数、谱与 Fourier 结构''}与\hyperref[ch:025]{``张量：多模结构怎样保留与压缩''}适合按需回补：前者要到信号、频域方法和振荡模态处才显出必要性，后者在处理多路数据时再看不迟。

读这一部分值得养成一个习惯：每遇到一个分解，先问它要求矩阵满足什么、换到手的新基好在哪里、以及要求不满足时退化成什么。九章的分解定理摆在一起，比较的正是这三件事。
